% Generated draft: requires mathematical and source review.
\documentclass[11pt]{article}
\usepackage[margin=25mm]{geometry}
\usepackage{fontspec}
\setmainfont{Latin Modern Roman}
\usepackage{amsmath,amssymb,mathtools,unicode-math}
\setmathfont{Latin Modern Math}
\usepackage{enumitem,needspace,xurl,hyperref}
\setlength{\parindent}{0pt}
\setlength{\parskip}{5pt}
\setlength{\emergencystretch}{3em}
\newcommand{\N}{\mathbb{N}}
\newcommand{\Z}{\mathbb{Z}}
\newcommand{\Q}{\mathbb{Q}}
\newcommand{\R}{\mathbb{R}}
\newcommand{\C}{\mathbb{C}}
\newcommand{\F}{\mathbb{F}}
\newcommand{\A}{\mathbb{A}}
\newcommand{\PP}{\mathbb{P}}
\newcommand{\E}{\mathbb{E}}
\newcommand{\eps}{\varepsilon}
\newcommand{\dd}{\,\mathrm{d}}
\begin{document}
\section*{UnsolvedMath ID: 30005372}
\textbf{Source problem code:} catalogue entry
\subsection{Definitions and mathematical statement}
Let \(X\) be a closed, simply connected, smooth \(4\)-manifold with a spin structure; in particular, \(X\) is oriented and \(w_2(TX)=0\). Here closed means compact without boundary, and simply connected includes connectedness.

The supplied question asks which such \(X\) have boundary twist \(t_X\) smoothly isotopic to the identity. However, the supplied excerpt does not define \(t_X\). Since \(X\) itself has no boundary, the manifold on which this twist acts and the boundary conditions for the isotopy must be specified.

A standard candidate interpretation, stated here explicitly as a reconstruction requiring source confirmation, is the following. Choose a smoothly embedded closed ball \(D^4\subset X\) and set
\[
X^{\circ}=X\setminus\operatorname{int}(D^4),\qquad \partial X^{\circ}\cong S^3.
\]
Choose a collar \(S^3\times[0,1]\subset X^{\circ}\), with \(S^3\times\{0\}=\partial X^{\circ}\). Let \(\gamma:[0,1]\to\operatorname{SO}(4)\) be a smooth based loop, constant at the identity near both endpoints, representing the nonzero element of \(\pi_1(\operatorname{SO}(4))\cong\Z/2\Z\). Using the standard action on \(S^3\subset\R^4\), define a diffeomorphism by
\[
t_X(u,s)=(\gamma(s)u,s)
\]
on the collar and by the identity elsewhere. This diffeomorphism fixes the boundary pointwise.

Under this reconstructed interpretation, the classification question is: for which closed, spin, simply connected, smooth \(4\)-manifolds \(X\) does there exist a smooth isotopy \(F:X^{\circ}\times[0,1]\to X^{\circ}\) such that, writing \(F_r=F(\cdot,r)\),
\[
F_0=t_X,\qquad F_1=\operatorname{id}_{X^{\circ}},\qquad
F_r\text{ is a diffeomorphism},\qquad
F_r|_{\partial X^{\circ}}=\operatorname{id}_{\partial X^{\circ}}
\quad\text{for every }r\in[0,1]?
\]
The source must confirm both this construction and the requirement that the isotopy fix the boundary pointwise. It must also clarify the intended dependence on the choices used in constructing the twist. This conditional formulation does not resolve those omissions. The adjacent discussion of homeomorphisms with unrestricted boundary behavior is not sufficient to determine the intended smooth isotopy convention.
\subsection{Short English statement}
The question asks which closed, spin, simply connected, smooth four-manifolds have a smoothly trivial boundary twist. A conventional interpretation uses a twist on the manifold with a ball removed and an isotopy fixing its boundary pointwise. The supplied excerpt does not define the twist or specify that isotopy convention, so this interpretation requires confirmation. The supplied record labels the general problem open; this draft does not independently verify that status.
\subsection{Sources}
\begin{itemize}
\item Local registry: \texttt{data/unsolvedmath/problems.json}, with \texttt{catalog/problems.json} as the additional project registry; supplied permanent UnsolvedMath ID \(30005372\), source problem code \texttt{catalogue entry}, Question 3. These registry references identify the supplied record; their files were not independently inspected.
\item Supplied primary citation: \emph{Low-dimensional Topology} (2023), Oberwolfach Reports, DOI 10.4171/owr/2023/3, pp.\ 151--213. Original Oberwolfach report: \url{https://doi.org/10.4171/owr/2023/3}. No individual proposer is identified in the supplied record.
\item Supplied literature source: Powell, mapping class groups of simply connected \(4\)-manifolds, \url{https://www.maths.gla.ac.uk/~mpowell/pseudoisotopy.pdf}. The record does not provide complete bibliographic details.
\item Supplied literature source: Lin, Dehn twist on K3 connected sums, \url{https://doi.org/10.2140/gt.2023.27.1987}.
\item Supplied contextual attribution: the excerpt credits Giansiracusa [Gia08] with an explicit geometric argument for smooth triviality of \(t_{\mathbb{CP}^2}\). It states that an extension to all non-spin, smooth, simply connected, closed \(4\)-manifolds was communicated by Auckly, Kronheimer, and Ruberman. It attributes nontriviality of \(t_{K3}\) independently to Baraglia--Konno [BK22] and Kronheimer--Mrowka [KM20]. The supplied record does not give full bibliographic entries for these reference labels; these claims and attributions are retained as source assertions, not independently verified results.
\item Literature assessment dated exactly \texttt{2026-08-22}: supplied dated curation, not independently verified. It labels the general spin question open and reports smooth triviality in the non-spin case and nontriviality for K3 and related connected sums. The assessment explicitly calls for independent verification of time-sensitive status and cited claims; it does not specify which related connected sums are covered.
\end{itemize}
\end{document}
